\documentclass[11pt,a4paper]{article}
\usepackage[T1]{fontenc}
\usepackage[utf8]{inputenc}
\usepackage[english]{babel}
\usepackage{amsmath,amssymb,longtable,array}
\usepackage[margin=26mm]{geometry}
\usepackage[colorlinks=true,urlcolor=blue,linkcolor=blue]{hyperref}
\title{Supplement S1\\Partial Lean Verification}
\author{Andrea Comparini\\\small A curious amateur with proof obligations}
\date{9 October 2026}
\begin{document}
\maketitle

\noindent Accompanying \emph{Covariogram Uniqueness with One Regular Body in
$\mathbb R^3$: A Candidate Proof via Area Measures and Fourier--Laplace Zeros}.
This supplement and the Lean sources are separate from the mathematical paper.

\section{What the verification establishes}
The project contains 33 explicitly stated theorems in nine modules.
All have been elaborated and checked by Lean 4.24.0, with warnings treated
as errors. The transitive axiom audit covers 71 theorem declarations when
automatically generated helpers are included. Only
\texttt{propext}, \texttt{Classical.choice}, and \texttt{Quot.sound}
occur. No admitted proof or added mathematical axiom is used.
The source checksums recorded with the successful run were rechecked
while preparing the paper and still match.

This certifies the statements with their declared hypotheses.
It does not derive those hypotheses from equality of covariograms.
In particular, the project does not yet contain a formal statement of the
complete geometric uniqueness theorem, and neither central proposition
of the paper has been fully formalized.

\section{Correspondence with the mathematical paper}
\begin{longtable}{@{}p{0.24\textwidth}p{0.07\textwidth}p{0.63\textwidth}@{}}
\textbf{Module} & \textbf{Count} & \textbf{Checked content and scope}\\ \hline
\endfirsthead
\textbf{Module} & \textbf{Count} & \textbf{Checked content and scope}\\ \hline
\endhead
\path{Algebra} & 7 &
Unordered pairs from sums and products, logarithmic ratios, choice of order,
and the coefficient of the finite-difference kernel.
Related to the end of Appendix A, Appendix C.1, and Section 4.\\[5pt]
\path{AnalyticChoice} & 1 &
The analytic identity argument on an open preconnected subset of a real
normed space. Applying it to the spherical branches in Appendix B.1
remains a separate task.\\[5pt]
\path{ZeroAllocation} & 2 &
Allocation of a simple product zero to exactly one factor, and exclusivity
of the global analytic choice. The differentiation and simplicity data
are hypotheses.\\[5pt]
\path{LogKernel} & 6 &
Normalization of the logarithmic kernel, a cubic remainder estimate,
the denominator correction, and the bound with constant 36 in Appendix C.1.\\[5pt]
\path{KernelSeries} & 5 &
Summability and vanishing of the total remainder; the splitting of the
convergent logarithmic and rational series. Bounded heights and summable
inverse squares of real parts are explicit hypotheses.\\[5pt]
\path{RealParts} & 3 &
Replacement of real parts by a model at bounded distance, using a summable
majorant and dominated convergence. Related to Appendix C.3.\\[5pt]
\path{WeightedErrors} & 3 &
Summability of bounded weighted sequences, the signed-mass bound, and the
finite-head/tail limit argument of Appendix C.2 for abstract positive weights.\\[5pt]
\path{LimitTransport} & 4 &
A common scalar from prescribed directional limits; its modulus bound;
the passage from local signs to one orientation almost everywhere.
The directional asymptotic identity is assumed.\\[5pt]
\path{MeasureConclusion} & 2 &
Orientation of densities from sum, product and transport identities, and
equality of the measures defined by these densities. Their identification
with surface area measures and Minkowski uniqueness are not formalized.\\ \hline
\end{longtable}

The main paper uses a finite difference to remove the exponential
coefficient in a Hadamard product. The checked algebraic kernel identity,
logarithmic linearization and general limit arguments support this route;
they do not constitute a formal proof of the entire finite-difference
construction or of the normalized mean in Appendix D.

\section{Representative checked statements}
For real $x\ne0$, $|y|\le M$, $\tau>0$ and $\tau\ge4M$, Lean verifies
\[
 \left|\log\frac{x^2+(\tau-y)^2}{x^2+(\tau+y)^2}
             +\frac{4\tau y}{x^2+\tau^2}\right|
 \le\frac{36M^3\tau}{(x^2+\tau^2)^2}.
\]
For $x_n\ne0$, $|y_n|\le M$ and $\sum_nx_n^{-2}<\infty$, it verifies
\[
 \sum_n\log\frac{x_n^2+(\tau-y_n)^2}{x_n^2+(\tau+y_n)^2}
       +4\tau\sum_n\frac{y_n}{x_n^2+\tau^2}\longrightarrow0.
\]
Both convergence of the needed series and the splitting used in this
statement are checked.

Another result concerns a full-support measure $\mu$ on a topological
space and real functions $a,b,p,q,w$. With $a,b>0$ and $w\ne0$, put
\[
 \Gamma(a,b,w)=\frac{\log a-\log b}{2w}.
\]
Assume that $\Gamma(a,b,w)$ is continuous and that, almost everywhere,
\[
 p+q=a+b,\qquad pq=ab,\qquad
 \Gamma(p,q,w)=c\,\Gamma(a,b,w)
\]
for one scalar $c$. Lean proves that either $(p,q)=(a,b)$ almost
everywhere or $(p,q)=(b,a)$ almost everywhere. It then identifies the
corresponding measures defined by densities. This result does not assert
that an arbitrary convex competitor satisfies the three hypotheses.

\section{Theorem inventory}
All names below belong to the namespace \texttt{CovariogramCheck}.
The files are in \path{formalizzazione/CovariogramCheck/}.

\subsection*{Algebra.lean}
{\small\begin{verbatim}
pair_from_sum_product
gamma_swap
gamma_eq_zero_iff
ordered_pair_from_gamma
reversed_pair_from_gamma
scalar_is_sign
finite_difference_kernel
\end{verbatim}}

\subsection*{AnalyticChoice.lean and ZeroAllocation.lean}
{\small\begin{verbatim}
analytic_factor_choice
simple_product_root_allocation
analytic_factor_choice_exclusive
\end{verbatim}}

\subsection*{LogKernel.lean}
{\small\begin{verbatim}
log_ratio_cubic_remainder
log_ratio_cubic_remainder_half
logKernel_normalization
logKernel_cubic_bound
kernel_denominator_correction
logKernel_uniform_remainder
\end{verbatim}}

\subsection*{KernelSeries.lean}
{\small\begin{verbatim}
logKernel_remainder_majorant
logKernel_remainder_sum_bound
logKernel_remainder_sum_tendsto_zero
linearKernel_summable
logKernel_sum_linearization
\end{verbatim}}

\subsection*{RealParts.lean}
{\small\begin{verbatim}
rationalKernel_difference_bound
rationalKernel_tendsto_zero
realParts_replacement_tendsto_zero
\end{verbatim}}

\subsection*{WeightedErrors.lean}
{\small\begin{verbatim}
summable_weighted_of_bounded
signed_sum_bound
weighted_errors_tendsto_zero
\end{verbatim}}

\subsection*{LimitTransport.lean}
{\small\begin{verbatim}
common_scalar_of_limits
limit_abs_le_one
scalar_is_sign_ae
density_orientation_ae
\end{verbatim}}

\subsection*{MeasureConclusion.lean}
{\small\begin{verbatim}
density_orientation_of_sum_product_transport
measure_orientation_of_sum_product_transport
\end{verbatim}}

\section{Unformalized dependencies}
The following links remain outside the certification:
\begin{enumerate}
\item The geometric definitions and the derivation of difference-body,
      brightness and Fourier--Laplace identities from equal covariograms.
\item The support-regularity argument, the area formula as applied to
      the whole surface area measure, and the section asymptotics in
      Appendix A, including the singular quadratic case.
\item The spectral results cited from Bianchi and their application to
      the analytic branches on the sphere, including exceptional zeros.
\item Hadamard factorization, the paired infinite product, and the
      logarithmic identity for the actual transforms in Appendix B.
\item The endpoint Laplace limits and their use to obtain positive
      coefficients and products of area densities.
\item The application of the abstract series estimates to all the
      parameters in the paper, and normalization of the specific weights
      in Appendix D.
\item Identification of the density-defined measures with geometric
      surface area measures, followed by Minkowski uniqueness and the
      centrally symmetric case.
\end{enumerate}
Some underlying tools exist in mathlib, but their existence in the library
does not amount to a formalized application to this problem.

\section{Versions, records and repetition}
The checked environment uses:
\begin{itemize}
\item Lean 4.24.0, commit
      \path{797c613eb9b6d4ec95db23e3e00af9ac6657f24b};
\item mathlib v4.24.0, commit
      \path{f897ebcf72cd16f89ab4577d0c826cd14afaafc7}.
\end{itemize}
The initial successful run was recorded at 23:03:48 UTC on 8 October 2026
(9 October in the author's local time). The current verification records
include a subsequent run of the publication checking script; its exact
time and environment are recorded in \path{versions.txt}.
The directory \path{formalizzazione/verifica/} contains
\path{build.log}, \path{kernel.log}, \path{axioms.log},
\path{sources.sha256} and \path{versions.txt}.

On the prepared local machine, running \path{formalizzazione/check.sh}
builds the project, directly rechecks every source with warnings as errors,
audits the transitive axioms, and writes the version and checksum records.
The audit rejects any axiom other than the three named above.
The option \texttt{autoImplicit = false} prevents undeclared identifiers
from silently introducing extra implicit parameters.

The current Lake configuration points to a local mathlib checkout in
\path{.tools/mathlib4/} relative to the repository root. The checking script
accepts Lean and Lake on the command path, or a binary directory specified
by \texttt{LEAN\_BIN}, with a fallback for the original macOS arm64 setup.
It checks the Lean version and mathlib revision before building.
On another computer, these dependencies must first be provided at the
pinned versions; the repository README gives the setup commands.
With the matching Lean and Lake executables on the command path and
mathlib in the configured location, the essential commands from
\path{formalizzazione/} are:
{\small\begin{verbatim}
lake build
lake env lean -DautoImplicit=false -DwarningAsError=true Audit.lean
\end{verbatim}}
Individual sources can be rechecked using the same \texttt{lake env lean}
command and their filenames. A fresh-machine installation has not been
tested as part of this manuscript preparation. No portability claim or
certification of the whole candidate theorem is inferred from the local run.
\end{document}
